\section{Esecuzione In-Memory con
Mascheramento}\label{esecuzione-in-memory-con-mascheramento}

Un'altra alternativa per evitare il costo computazionale e di storage
della clonazione dei container consiste nell'eseguire gli script di
verifica direttamente all'interno dell'ambiente di produzione,
mascherando l'esecuzione su piu fronti, estendendo l'idea vista in 4.3
per creare un'ambiente completamente separato per il payload.

\textbf{esecuzione in memoria}: Si mette lo script o il binario in RAM,
utilizzando \passthrough{\lstinline!memfd\_create!}, e lo si esegue
usando \passthrough{\lstinline!fexecve!}, si evitano file persistenti
sul filesystem.

\textbf{mascheramento di \passthrough{\lstinline!/proc!}}: per
nascondere il processo in esecuzione a \passthrough{\lstinline!pspy!} si
inietta nel sistema una libreria a livello utente
\passthrough{\lstinline!/etc/ld.so.preload!} oppure modificando la
variabile d'ambiente \passthrough{\lstinline!LD\_PRELOAD!},
intercettando le funzioni della lib \passthrough{\lstinline!C!} che
leggono dentro \passthrough{\lstinline!/proc!}.

\textbf{mascheramento dei metadati del processo}: sovrascrivere il nome
del processo nelle strutture dati del kernel ancora prima che venga
avviato e catturato da \passthrough{\lstinline!pspy!}.

tutta via, queste sono tecniche di oscuramento dell'esecuzione e possono
essere facilmente mascherate e potrebbero risultare incompatibili con i
requisiti dell'utente:

\begin{itemize}
\tightlist
\item
  l'exploit potrebbe totalmente ignorare
  \passthrough{\lstinline!LD\_PRELOAD!}.
\item
  il file \passthrough{\lstinline!/etc/ld.so.preload!} potrebbe essere
  accessibile allo sfidante, e dunque rimovibile.
\end{itemize}

a livello kernel, esiste comunque un traccia dell'esecuzione del
payload: se l'utente e' in grado di registrare l'esecuzione delle
chiamate di sistema allora puo' catturare l'esecuzione del payload.

Anche se l'eseguibile del verificatore risiede in RAM, se questo inoltra
richiesta HTTP con all'interno il payload, allora il contenuto puo'
essere catturato e letto dallo sfidante.

\begin{itemize}
\tightlist
\item
  il traffico generato puo' essere letto da
  \passthrough{\lstinline!tcpdump!}
\item
  oppure la richiesta HTTP puo' modificare file accessibili allo
  sfidante.
\end{itemize}

l'esecuzione dell'exploit in ambiente separato potrebbe comunque portare
in crash il servizio.

dunque, l'esecuzione in-memory con mascheramento di
\passthrough{\lstinline!/proc!} si basa su tecniche di evasione fragili
(``security through obscurity''). Non offre un vero isolamento: non
protegge i file di log, non isola il traffico di rete e non previene il
crash distruttivo dei servizi live.

La clonazione tramite ZFS CoW in una Dark Sandbox OVN rimane la
soluzione superiore, poiché garantisce un isolamento deterministico a
livello di kernel e hypervisor senza ricorrere a trucchi di evasione in
user-space.
